% introduccion.tex
\section{Introducción}

\subsection{Contexto del Problema}
Los sensores de imagen en cámaras digitales utilizan filtros de color para capturar información cromática. 
El patrón Bayer, desarrollado por Bryce Bayer en 1976, es el arreglo más común, donde cada fotosito del sensor 
captura únicamente un canal de color (rojo, verde o azul).

\subsection{Problema de Interpolación}
Dada una matriz 2×2 del patrón Bayer:

\[
\begin{bmatrix}
R & G \\
G & B
\end{bmatrix}
\]

Es necesario estimar los valores faltantes en cada posición para reconstruir tres canales completos (RGB) 
con la misma resolución espacial. Este proceso se denomina \textbf{demosaicing} o interpolación de Bayer.

\subsection{Objetivos del Estudio}
\begin{enumerate}
    \item Implementar algoritmos de interpolación bicuadrada para el patrón Bayer RGGB
    \item Implementar algoritmos de interpolación bicúbica para el patrón Bayer RGGB
    \item Comparar los resultados con métodos de interpolación lineal tradicional
    \item Analizar la complejidad computacional de cada método
    \item Visualizar los procesos mediante diagramas y gráficos explicativos
    \item Proporcionar implementaciones prácticas en MATLAB para uso educativo
\end{enumerate}

\subsection{Organización del Documento}
Este documento se estructura de la siguiente manera: la Sección \ref{sec:metodos} describe los métodos 
matemáticos utilizados; la Sección \ref{sec:resultados} presenta los resultados obtenidos; la Sección 
\ref{sec:discusion} analiza los hallazgos; y finalmente, la Sección \ref{sec:conclusion} ofrece conclusiones 
y trabajo futuro. Los archivos de código fuente se encuentran en el Apéndice.